\section{2013级工科数学分析下 B卷}

\subsection{资料来源}
\begin{itemize}
  \item 试题：2013 级软件工科数学分析下 B 卷，已导出页面图校正公式。
  \item 未发现配套答案或解析文本；下列答案与解析为据题面补写。
\end{itemize}

\subsection{试题}

\paragraph*{试卷信息}
诚信应考，考试作弊将带来严重后果！

华南理工大学本科生期末考试《工科数学分析》2013--2014学年第二学期期末考试试卷（B）卷。

注意事项：开考前请将密封线内各项信息填写清楚；所有答案请直接答在试卷上（或答题纸上）；考试形式：闭卷；本试卷共5个大题，满分100分，考试时间120分钟。

\paragraph*{一、单项选择题（每小题3分，共15分）}
\begin{enumerate}
  \item 方程 \((y-\ln x)\,\mathrm{d}x+x\,\mathrm{d}y=0\) 是（\quad）。
  \begin{enumerate}
    \item[A.] 可分离变量微分方程。
    \item[B.] 一阶线性非齐次方程。
    \item[C.] 一阶线性齐次方程。
    \item[D.] 非线性方程。
  \end{enumerate}

  \item 设 \(f(x,y)\) 是可微函数，且 \(f(0,0)=0\)，\(f_x(0,0)=a\)，\(f_y(0,0)=b\)，令
  \(\varphi(t)=f(t,f(t,t))\)，则 \(\varphi'(0)=\)（\quad）。
  \begin{enumerate}
    \item[A.] \(a\)
    \item[B.] \(a+b(a+b)\)
    \item[C.] \(a+1\)
    \item[D.] \(\dfrac{a}{1-b}\)
  \end{enumerate}

  \item 设 \(C:x^2+y^2=a^2\)，取逆时针方向，则
  \(\displaystyle \oint_C xy^2\,\mathrm{d}y-x^2y\,\mathrm{d}x=\)（\quad）。
  \begin{enumerate}
    \item[A.] \(\dfrac{\pi a^4}{2}\)
    \item[B.] \(-\dfrac{\pi a^4}{2}\)
    \item[C.] \(\pi a^4\)
    \item[D.] \(-\pi a^4\)
  \end{enumerate}

  \item 幂级数 \(\displaystyle \sum_{n=1}^{\infty}\frac{x^n}{n2^n}\) 的收敛域是（\quad）。
  \begin{enumerate}
    \item[A.] \((-2,2)\)
    \item[B.] \([-2,2]\)
    \item[C.] \((-2,2]\)
    \item[D.] \([-2,2)\)
  \end{enumerate}

  \item 设 \(f(x)\) 是周期为 \(2\) 的周期函数，它在区间 \((-1,1]\) 上的表达式为
  \[
    f(x)=
    \begin{cases}
      2, & -1<x\le 0,\\
      x^3, & 0<x\le 1,
    \end{cases}
  \]
  则 \(f(x)\) 的 Fourier 级数在 \(x=1\) 处收敛于（\quad）。
  \begin{enumerate}
    \item[A.] \(1\)
    \item[B.] \(0\)
    \item[C.] \(\dfrac32\)
    \item[D.] \(-\dfrac32\)
  \end{enumerate}
\end{enumerate}

\paragraph*{二、填空题（每小题3分，共15分）}
\begin{enumerate}
  \item 微分方程 \(y''-2y'+y=0\) 的通解为 \underline{\hspace{5cm}}。
  \item 函数 \(\displaystyle z=\arctan\frac{y}{x}+\ln\sqrt{x^2+y^2}\) 在点 \((1,2)\) 处的梯度为 \underline{\hspace{5cm}}。
  \item 交换积分的顺序，则 \(\displaystyle \int_0^2\mathrm{d}x\int_x^2 \mathrm{e}^{-y^2}\,\mathrm{d}y=\) \underline{\hspace{5cm}}。
  \item 设 \(\Sigma:x^2+y^2+z^2=a^2\)，则
  \(\displaystyle \iint_\Sigma (x^2+y^2+z^2)\,\mathrm{d}S=\) \underline{\hspace{5cm}}。
  \item 设幂级数
  \[
    \sum_{n=1}^{\infty}a_n(x-1)^n
  \]
  在 \(x=-1\) 处收敛，则此级数在 \(x=2\) 的敛散性是 \underline{\hspace{5cm}}。
\end{enumerate}

\paragraph*{三、计算题（每小题10分，共50分）}
\begin{enumerate}
  \item 设 \(f(\xi,\eta)\) 具有连续的二阶偏导数，且满足 Laplace 方程
  \(\displaystyle \frac{\partial^2 f}{\partial \xi^2}+\frac{\partial^2 f}{\partial \eta^2}=0\)，令
  \(z=f(x^2-y^2,2xy)\)，求
  \(\displaystyle \frac{\partial^2z}{\partial x^2}+\frac{\partial^2z}{\partial y^2}\)。

  \item 计算三重积分
  \[
    I=\iiint_\Omega z\mathrm{e}^{-(x^2+y^2+z^2)}\,\mathrm{d}V,
    \qquad
    \Omega:\begin{cases}
      1\le x^2+y^2+z^2\le 4,\\
      x\ge 0,\ y\ge 0,\ z\ge 0.
    \end{cases}
  \]

  \item 已知 \(\varphi(\pi)=1\)，试确定 \(\varphi(x)\)，使曲线积分
  \[
    I=\int_A^B\left[\sin x-\varphi(x)\right]\frac{y}{x}\,\mathrm{d}x+\varphi(x)\,\mathrm{d}y
  \]
  与路径无关，并求当 \(A,B\) 两点分别为 \((1,0),(\pi,\pi)\) 时积分 \(I\) 的值。

  \item 计算曲面积分
  \[
    I=\iint_\Sigma (8y+1)x\,\mathrm{d}y\mathrm{d}z+2(1-y^2)\,\mathrm{d}z\mathrm{d}x-4yz\,\mathrm{d}x\mathrm{d}y,
  \]
  其中 \(\Sigma:y=1+x^2+z^2\)，介于 \(1\le y\le 3\) 之间，且其法向量与 \(y\) 轴正向的夹角恒大于 \(\dfrac{\pi}{2}\)。

  \item 求级数 \(\displaystyle \sum_{n=1}^{\infty}\frac{(-1)^n n}{(2n+1)!}\) 的和。
\end{enumerate}

\paragraph*{四、证明题（本题10分）}
\begin{enumerate}
  \setcounter{enumi}{5}
  \item 证明级数 \(\displaystyle \sum_{n=1}^{\infty}\frac{\sin(nx)}{n^2}\) 在 \((-\infty,+\infty)\) 上一致收敛。
\end{enumerate}

\paragraph*{五、应用题（本题10分）}
\begin{enumerate}
  \setcounter{enumi}{6}
  \item 求内接于半径为 \(a\) 的球且有最大体积的长方体。
\end{enumerate}

\subsection{答案与解析}
\paragraph*{一、单项选择题}
\begin{enumerate}
  \item 选 B。原方程可化为
  \[
    y'+\frac1x y=\frac{\ln x}{x},
  \]
  因而是一阶线性非齐次方程。

  \item 选 B。由链式法则，
  \[
    \varphi'(t)=f_x(t,f(t,t))+f_y(t,f(t,t))\frac{\mathrm{d}}{\mathrm{d}t}f(t,t),
  \]
  且 \(\left.\dfrac{\mathrm{d}}{\mathrm{d}t}f(t,t)\right|_{t=0}=f_x(0,0)+f_y(0,0)=a+b\)，所以
  \(\varphi'(0)=a+b(a+b)\)。

  \item 选 A。令 \(P=-x^2y\)，\(Q=xy^2\)。由 Green 公式，
  \[
    \oint_C P\,\mathrm{d}x+Q\,\mathrm{d}y
    =\iint_D(Q_x-P_y)\,\mathrm{d}A
    =\iint_D(x^2+y^2)\,\mathrm{d}A
    =\frac{\pi a^4}{2}.
  \]

  \item 选 D。由比值判别法得收敛半径 \(R=2\)。当 \(x=2\) 时为调和级数 \(\sum 1/n\)，发散；当 \(x=-2\) 时为交错调和级数，收敛。因此收敛域为 \([-2,2)\)。

  \item 选 C。周期为 \(2\)，在 \(x=1\) 处按 Fourier 收敛定理取左右极限平均值：
  \[
    \frac{f(1-0)+f(1+0)}2=\frac{1+2}{2}=\frac32.
  \]
\end{enumerate}

\paragraph*{二、填空题}
\begin{enumerate}
  \item 特征方程 \((r-1)^2=0\)，故通解为
  \[
    y=(C_1+C_2x)\mathrm{e}^x.
  \]

  \item 有
  \[
    z_x=-\frac{y}{x^2+y^2}+\frac{x}{x^2+y^2},\qquad
    z_y=\frac{x}{x^2+y^2}+\frac{y}{x^2+y^2}.
  \]
  所以
  \[
    \nabla z(1,2)=\left(-\frac15,\frac35\right).
  \]

  \item 积分区域为 \(0\le x\le y\le 2\)，故
  \[
    \int_0^2\mathrm{d}x\int_x^2 \mathrm{e}^{-y^2}\,\mathrm{d}y
    =\int_0^2\mathrm{d}y\int_0^y \mathrm{e}^{-y^2}\,\mathrm{d}x
    =\int_0^2y\mathrm{e}^{-y^2}\,\mathrm{d}y
    =\frac{1-\mathrm{e}^{-4}}2.
  \]

  \item 在球面上 \(x^2+y^2+z^2=a^2\)，故
  \[
    \iint_\Sigma (x^2+y^2+z^2)\,\mathrm{d}S=a^2\cdot 4\pi a^2=4\pi a^4.
  \]

  \item 该幂级数以 \(x=1\) 为中心，在 \(x=-1\) 处收敛，说明收敛半径至少为 \(2\)。点 \(x=2\) 到中心的距离为 \(1<2\)，故级数在 \(x=2\) 处绝对收敛。
\end{enumerate}

\paragraph*{三、计算题}
\begin{enumerate}
  \item 记 \(\xi=x^2-y^2,\eta=2xy\)。由链式法则直接计算可得
  \[
    \frac{\partial^2 z}{\partial x^2}+\frac{\partial^2 z}{\partial y^2}
    =4(x^2+y^2)\left(\frac{\partial^2 f}{\partial \xi^2}+\frac{\partial^2 f}{\partial \eta^2}\right)=0.
  \]

  \item 用球坐标
  \(x=r\sin\varphi\cos\theta,\ y=r\sin\varphi\sin\theta,\ z=r\cos\varphi\)。区域为
  \(1\le r\le 2,\ 0\le\varphi\le\pi/2,\ 0\le\theta\le\pi/2\)，因此
  \[
    I=\int_0^{\pi/2}\!\!\int_0^{\pi/2}\!\!\int_1^2
    r\cos\varphi\,\mathrm{e}^{-r^2}r^2\sin\varphi\,\mathrm{d}r\,\mathrm{d}\varphi\,\mathrm{d}\theta
    =\frac{\pi}{8}\left(\frac2{\mathrm{e}}-\frac5{\mathrm{e}^4}\right).
  \]

  \item 设
  \[
    P=\left[\sin x-\varphi(x)\right]\frac{y}{x},\qquad Q=\varphi(x).
  \]
  路径无关要求 \(P_y=Q_x\)，即
  \[
    \frac{\sin x-\varphi(x)}{x}=\varphi'(x),
    \qquad (x\varphi(x))'=\sin x.
  \]
  积分得 \(x\varphi(x)=-\cos x+C\)。由 \(\varphi(\pi)=1\) 得 \(C=\pi-1\)，所以
  \[
    \varphi(x)=\frac{\pi-1-\cos x}{x}.
  \]
  势函数可取 \(F(x,y)=y\varphi(x)\)，故
  \[
    I=F(\pi,\pi)-F(1,0)=\pi.
  \]

  \item 取参数 \((x,z)\)，令 \(y=1+x^2+z^2\)。题设法向量的 \(y\) 分量为负，故取
  \[
    \mathbf r_x\times\mathbf r_z=(2x,-1,2z).
  \]
  投影区域为 \(D:x^2+z^2\le 2\)。代入第二类曲面积分，
  \[
    I=\iint_D\left[(8y+1)x\cdot2x-2(1-y^2)-4yz\cdot2z\right]\mathrm{d}x\mathrm{d}z,
    \qquad y=1+x^2+z^2.
  \]
  化简后利用圆盘对称性：
  \[
    I=10\iint_Dx^2\,\mathrm{d}A+4\iint_Dr^2\,\mathrm{d}A+6\iint_Dr^4\,\mathrm{d}A
    =10\pi+8\pi+16\pi=34\pi.
  \]

  \item 由
  \[
    \sin t=\sum_{n=0}^{\infty}\frac{(-1)^n t^{2n+1}}{(2n+1)!}
  \]
  得
  \[
    \sum_{n=1}^{\infty}\frac{(-1)^n n}{(2n+1)!}
    =\frac12\sum_{n=1}^{\infty}\frac{(-1)^n(2n+1)}{(2n+1)!}
     -\frac12\sum_{n=1}^{\infty}\frac{(-1)^n}{(2n+1)!}
    =\frac{\cos1-\sin1}{2}.
  \]
\end{enumerate}

\paragraph*{四、证明题}
\begin{enumerate}
  \setcounter{enumi}{5}
  \item 对任意 \(x\in(-\infty,+\infty)\)，都有
  \[
    \left|\frac{\sin(nx)}{n^2}\right|\le \frac1{n^2}.
  \]
  又数项级数 \(\sum_{n=1}^{\infty}1/n^2\) 收敛，由 Weierstrass 判别法，级数
  \(\sum_{n=1}^{\infty}\sin(nx)/n^2\) 在 \((-\infty,+\infty)\) 上一致收敛。
\end{enumerate}

\paragraph*{五、应用题}
\begin{enumerate}
  \setcounter{enumi}{6}
  \item 设长方体以球心为中心，三个半边长为 \(x,y,z>0\)，则约束为
  \[
    x^2+y^2+z^2=a^2,
  \]
  体积为 \(V=8xyz\)。由对称性或 Lagrange 乘子法，最大值在
  \(x=y=z=a/\sqrt3\) 处取得。因此长方体三条棱长均为
  \[
    \frac{2a}{\sqrt3},
  \]
  最大体积为
  \[
    V_{\max}=8\left(\frac{a}{\sqrt3}\right)^3=\frac{8a^3}{3\sqrt3}.
  \]
\end{enumerate}
